\section{Introduction}
\begin{itemize}
    \item Graphs neural networks are used for ... 
    \item Recently, there has been increased interest in using graph neural networks for solving mathematical, and combinatorial problems, ... 

    \item Examples of successful results in this vein are ... 
    \item The clique-colouring principle is ... 
    \item This principle has been a fundamental tool in proof complexity theory, and it has been shown to be hard for ... (lots of proof systems. See my paper with Pavel Pudlak)
    \item It is worth noting that graph neural networks can compute non-monotone functions, such as ... 
    \item Graph neural networks can be generalised to r-dimensions, with increasing expressive power. 

\end{itemize}

\subsection{Random graphs: }

\begin{enumerate}
    \item The clique colouring principle is very easy for (Normal) neural networks on the average case under the stardard distribution, that means where we select an edge with probability $1/2$. 
         (Check whether this works also for graphs neural networks. )
         (Check whether this works for other sampling probabilities. This may be the case for any fixed p)
         
    \item There is a specific distribution where (Normal) neural networks cannot distinguish the graphs. That is confirmed empirically.
        
\end{enumerate}


\subsection{Pinpointing speficic hard families:} 

\begin{enumerate}
    \item We show that for each $k\geq3$ no graph neural network can determine whether an arbitrary graph has a k-clique or whether it is k-colorable. More precisely, for each fixed k, we exhibit an 
infinite family of pairs of graphs ${(G_s^k,H_s^k)}_{s\in N}$  satisfying the following properties for each $s\in N$. 
    \begin{enumerate}
        \item $G_s^k$ contains a k-clique and $H_s^k$ is $(k-1)$-colorable
        \item No graph neural network can disguise $G_s^k$ from $H_s^k$. In other words, such a network N either accepts both graphs or rejects both graphs.
    \end{enumerate}
    \item We generalise our result for higher dimensions. More specifically, for any dimension r, and any k, there is a family of pairs of graphs ${(G[r,k,s],H[r,k,s])}_{s\in N}$ such that,
    \begin{enumerate}
        \item $G[r,k,s]$ contains a $k$-clique and $H[r,k,s]$ is (k-1)-colorable 
        \item No r-dimensional graph neural network can dustinguish $G[r,k,s]$ from $H[r,k,s]$
    \end{enumerate}
\end{enumerate}

\subsection{Some comments}

\begin{itemize}
    \item To be confirmed: The $(k-1)$-coloring problem is NP-complete for each $k\geq3$. Nevertheless, there is a polynomial time algorithm that distinguishes graphs 
that have $k$-cliques from graphs that are $(k-1)$-colorable. The algorithm uses the Lovas Theta function. (CHECK THIS! It may be the case that this works for a speficic subclass of graphs)

    CHECKED: For every fixed $k \ge 3$, deciding whether a graph is $(k-1)$-colorable is NP-complete. However, the \emph{promise problem} of distinguishing graphs that contain a $k$-clique from graphs that are $(k-1)$-colorable can be solved in polynomial time using the Lovász $\vartheta$-function: by the sandwich theorem $\omega(G) \le \vartheta(\overline{G}) \le \chi(G)$, so under the promise we obtain a gap---$\vartheta(\overline{G}) \ge k$ in the first case and $\vartheta(\overline{G}) \le k-1$ in the second---allowing us to separate the two via an efficiently computable approximation of $\vartheta$. This does not contradict NP-completeness because only the gap/promise version becomes tractable.

    \item Therefore, we show a class of graphs that that can be distinguished easily by polynomial time algorithms but are provably impossible from being distihguished by GNNS even in higher dimensions.

  \item  Our impossibility results are unconditional and do not require any complexity theoretic assumption. 
     
  
\end{itemize}



